(i) Each map below is an algebra homomorphism because the two factors act by commuting operators, and it is an isomorphism because its images span a space of the target's dimension, an exact rank computation recorded in the check. $\HH \otimes \HH \to \operatorname{End}_\R(\HH)$, $a \otimes b \mapsto (x \mapsto a x \bar{b})$: the $16$ maps on basis elements have rank $16$, and $\operatorname{End}_\R(\HH) = \Mat_4(\R)$. $\C \otimes \HH \to \operatorname{End}_\C(\HH)$, $z \otimes a \mapsto (x \mapsto a x z)$ with $\HH$ a right $\C$-space: rank $8 = \dim_\R \Mat_2(\C)$. $\C \otimes \C \to \C \oplus \C$, $z \otimes w \mapsto (zw, z\bar{w})$: rank $4$.

(ii) Base cases: $\Cl(0,0) = \R$; $\Cl(1,0) = \R e_+ \oplus \R e_-$ with the central orthogonal idempotents $e_{\pm} = \tfrac12(1 \pm e_1)$; $\Cl(0,1) \cong \C$ since $e_1^{2} = -1$; $\Cl(2,0) \cong \Mat_2(\R)$ by \claimref{std-algebra.cl11-matrix}; $\Cl(0,2) \cong \HH$ with $\mathrm{i} = e_1$, $\mathrm{j} = e_2$, $\mathrm{k} = e_1e_2$, which satisfy Hamilton's relations, the dimensions being $4$ (\claimref{std-algebra.blade-basis}). By \claimref{std-algebra.tensor-rules} with $V = \R^{0,k}$ or $\R^{k,0}$: $\Cl(k{+}2, 0) \cong \Cl(0,k) \otimes \Cl(2,0) = \Mat_2(\Cl(0,k))$ and $\Cl(0, k{+}2) \cong \Cl(k,0) \otimes \Cl(0,2) = \Cl(k,0) \otimes \HH$. Using (i) and $(A \oplus A) \otimes B = (A \otimes B) \oplus (A \otimes B)$, $\Mat_a(A) \otimes B = \Mat_a(A \otimes B)$, the two rows follow by induction on $k$: for example $\Cl(0,6) \cong \Cl(4,0) \otimes \HH = \Mat_2(\HH \otimes \HH) = \Mat_8(\R)$ and $\Cl(8,0) \cong \Mat_2(\Cl(0,6)) = \Mat_{16}(\R)$, $\Cl(0,8) \cong \Cl(6,0) \otimes \HH = \Mat_4(\HH \otimes \HH) = \Mat_{16}(\R)$.

(iii) The first isomorphism is \claimref{std-algebra.tensor-rules} with $W = \R^{1,1}$ and $\Cl(1,1) \cong \Mat_2(\R)$. The second is the same rule with $W = \R^{8,0}$ or $\R^{0,8}$, whose pseudoscalar squares to $+1$, and (ii). Applying the first $\min(p,q)$ times reduces $\Cl(p,q)$ to $\Mat_{2^{\min(p,q)}}(\Cl(p - q, 0))$ or $\Mat_{2^{\min(p,q)}}(\Cl(0, q - p))$, and the second reduces the remaining index below $8$; reading the table by $d = p - q \bmod 8$ ($\Cl(0,k)$ has $d = -k$) gives exactly the types of \claimref{primer.classification-real-clifford-algebras}.
